\documentclass[9pt,a4paper]{article}
\usepackage[french]{babel}
\usepackage[T1]{fontenc}
\usepackage{amsmath, amssymb}
\usepackage{geometry}
\geometry{top=1cm, bottom=2cm, left=1cm, right=1cm}
\usepackage{tcolorbox}
\usepackage{listings}
\usepackage{xcolor}

\definecolor{codegray}{rgb}{0.95,0.95,0.95}
\lstset{
    language=Python,
    backgroundcolor=\color{codegray},
    basicstyle=\ttfamily\small,
    keywordstyle=\color{blue},
    commentstyle=\color{gray},
    stringstyle=\color{red},
    frame=single,
    breaklines=true
}

\begin{document}

\centering\textbf{Feuille d’exercices d’entraînement - Algorithmique Python : Affectation et condition \texttt{if}}\\
\vspace{0.5cm}

\noindent\textbf{Nom :} \underline{\hspace{5cm}} \hfill \textbf{Prénom :} \underline{\hspace{5cm}} \\
\textbf{Classe :} \underline{\hspace{5cm}}\\
\vspace{1cm}

\section*{Exercice 1 — Lecture de programme}

\noindent Que va afficher le programme suivant après exécution ? Justifie ta réponse.

\begin{lstlisting}
temperature = 5

if temperature > 0:
    print("Au-dessus de zero")
else:
    print("En dessous de zero")
\end{lstlisting}

\begin{tcolorbox}[colframe=black, colback=white, width=\linewidth]
\vspace{3cm}
\end{tcolorbox}

\vspace{0.3cm}

\noindent Et si on remplace la première ligne par \texttt{temperature = 0} ? Que se passe-t-il ?

\begin{tcolorbox}[colframe=black, colback=white, width=\linewidth]
\vspace{3cm}
\end{tcolorbox}

\vspace{1cm}

\section*{Exercice 2 — Écriture de programme simple}

\noindent Écris un programme qui :
\begin{itemize}
    \item demande à l’utilisateur son âge ;
    \item affiche "\texttt{Majeur}" si l’âge est supérieur ou égal à 18 ;
    \item sinon affiche "\texttt{Mineur}".
\end{itemize}

\begin{tcolorbox}[colframe=black, colback=white, width=\linewidth]
\vspace{5cm}
\end{tcolorbox}

\vspace{1cm}

\section*{Exercice 3 — Conditions imbriquées}

\noindent Que va afficher le programme suivant ? Justifie.

\begin{lstlisting}
score = 17

if score >= 10:
    print("Admis")
    if score >= 16:
        print("Mention tres bien")
    else:
        print("Sans mention")
else:
    print("Refuse")
\end{lstlisting}

\begin{tcolorbox}[colframe=black, colback=white, width=\linewidth]
\vspace{4cm}
\end{tcolorbox}

\vspace{1cm}

\section*{Exercice 4 — Paires et impaires}

\noindent Écris un programme qui :
\begin{itemize}
    \item crée deux variables \texttt{a} et \texttt{b} avec les valeurs de ton choix ;
    \item calcule leur somme ;
    \item affiche "\texttt{La somme est paire}" si la somme est paire, sinon affiche "\texttt{La somme est impaire}".
\end{itemize}

\begin{tcolorbox}[colframe=black, colback=white, width=\linewidth]
\vspace{5cm}
\end{tcolorbox}

\vspace{0.3cm}
\textbf{Rappel :} Pour savoir si un nombre est pair, on utilise : \\
\texttt{if nombre \% 2 == 0:}

\vspace{1cm}

\newpage
\section*{Exercice 5 — Bonus (réflexion)}

\noindent Que va-t-il se passer avec ce programme ? Explique le résultat.

\begin{lstlisting}
x = 5
y = 10

if x + y > 10:
    print("Grande somme")
if x + y > 15:
    print("Tres grande somme")
else:
    print("Somme moyenne")
\end{lstlisting}

\begin{tcolorbox}[colframe=black, colback=white, width=\linewidth]
\vspace{4cm}
\end{tcolorbox}

\end{document}
